\documentclass[12pt]{report}
\usepackage[a4paper, left=3.17cm, right=3.17cm, top=2.54cm, bottom=2.54cm]{geometry}
\usepackage[T1]{fontenc}
\usepackage{mathptmx}
\usepackage{amsmath}
\usepackage{amsfonts}
\usepackage{multicol}
\usepackage{multirow}
\usepackage{tabularx,booktabs}
\newcolumntype{C}{>{\centering\arraybackslash}X}
\newcolumntype{P}[1]{>{\centering\arraybackslash}p{#1}}
\usepackage[linesnumbered,ruled,vlined]{algorithm2e}
\usepackage{comment}
\usepackage{array}
\usepackage{cite}
\usepackage[colorlinks, linkcolor=black, anchorcolor=black, citecolor=black]{hyperref}
\usepackage{graphicx}
\usepackage{xcolor}
\usepackage{listings}
\usepackage{enumitem}
\setlength{\parskip}{0.5em}

% Java code listing style
\lstdefinestyle{javaStyle}{
    language=Java,
    basicstyle=\ttfamily\small,
    keywordstyle=\color{blue}\bfseries,
    commentstyle=\color{gray}\itshape,
    stringstyle=\color{orange},
    numberstyle=\tiny\color{gray},
    numbers=left,
    stepnumber=1,
    frame=single,
    breaklines=true,
    captionpos=b,
    tabsize=4,
    showstringspaces=false
}

\title{Auralin Music Player}

% Copyright footer
\usepackage{fancyhdr}
\fancyhf{}
\rfoot{%
  \footnotesize
  \textcopyright~Dept. of Software Engineering, GUB\\
}

\begin{document}

    \input{title/title.tex}

    % Abstract
    \begin{abstract}
    \addcontentsline{toc}{chapter}{Abstract}
    Auralin Music Player is a feature-rich desktop music player application developed using Java and JavaFX. The application provides a modern, Spotify-inspired user interface with a dark theme, enabling users to manage and play their local MP3 music library with ease. The system implements a layered MVC (Model-View-Controller) architecture, separating concerns across dedicated packages for models, views, controllers, services, managers, components, and utilities. Key features include local file and folder import, background metadata extraction via the mp3agic library, a persistent library stored on disk, playback controls (play, pause, next, previous, shuffle, repeat), a liked songs collection, a home screen with recently played and recommended tracks, a full-screen Now Playing overlay, and a custom frameless window with native-style resize and drag support. This report documents the system architecture, design patterns, implementation details, database persistence strategy, UI design decisions, and the challenges encountered during development.
    \end{abstract}

    \tableofcontents


% ============================================================
% CHAPTER 1: INTRODUCTION
% ============================================================
\newpage
\chapter{Introduction}

\section{Overview}
Auralin Music Player is a desktop application built with Java 17 and JavaFX that allows users to manage and play their local MP3 music collection. The application features a sleek, dark-themed interface inspired by modern streaming platforms, offering a seamless music listening experience without requiring an internet connection. The player supports importing individual files or entire folders, automatically reads ID3 metadata tags (title, artist, album artwork), and persists the library between sessions.

\section{Motivation}
The primary motivation for developing Auralin was to create a lightweight, offline-capable music player that combines the aesthetic quality of modern streaming applications with the simplicity of a local media player. Existing open-source players often lack visual polish or are overly complex. Auralin aims to bridge this gap by providing:

\begin{itemize}
    \item A visually appealing, Spotify-inspired dark UI built entirely in JavaFX.
    \item Fast background metadata loading without blocking the UI thread.
    \item A persistent library that remembers songs, play counts, and liked status across sessions.
    \item A clean MVC architecture that demonstrates good object-oriented design principles.
    \item Support for the ``Open With'' shell integration so MP3 files can be opened directly from Windows Explorer.
\end{itemize}

\section{Problem Definition}

\subsection{Problem Statement}
Most existing Java-based music players either rely on outdated Swing UI toolkits, lack metadata support, or do not persist library state between sessions. Users who want a modern-looking, offline music player with a clean codebase face limited options. Auralin addresses these gaps by:

\begin{itemize}
    \item Providing a fully custom, frameless JavaFX window with rounded corners and native-style window controls.
    \item Loading ID3v1 and ID3v2 metadata (artist, artwork, duration) in background threads.
    \item Saving and restoring the library, play counts, liked status, and last-played timestamps to disk.
    \item Offering shuffle, repeat, and liked-songs filtering features.
\end{itemize}

\subsection{Complex Engineering Problem}
Table~\ref{tab:CEP} summarises the engineering attributes addressed by this project.

\begin{table}[htbp]
   \centering
    \caption{Summary of complex engineering problem attributes}
    \begin{tabular}{|p{5.5cm}|p{8.5cm}|}
    \toprule
        \textbf{Attribute} & \textbf{How it is addressed} \\
    \midrule
    \textbf{P1: Depth of knowledge} & Requires JavaFX scene graph, MediaPlayer API, ID3 tag parsing, multithreading, and file I/O. \\ \hline
    \textbf{P2: Conflicting requirements} & Balancing UI responsiveness with background metadata loading and disk persistence. \\ \hline
    \textbf{P3: Depth of analysis} & Analysing playback state transitions, observer callbacks, and thread-safe UI updates via \texttt{Platform.runLater()}. \\ \hline
    \textbf{P4: Familiarity of issues} & Addresses common issues: UI freezing during file scan, stale table rows after like-toggle, and artwork not loading. \\ \hline
    \textbf{P5: Applicable codes} & Applies MVC pattern, Observer pattern, Strategy pattern, and JavaFX property bindings. \\ \hline
    \textbf{P6: Stakeholder involvement} & End-users want simplicity; developers need extensibility --- resolved through clean interface boundaries. \\ \hline
    \textbf{P7: Interdependence} & Strong coupling between \texttt{PlayerController}, \texttt{DatabaseService}, \texttt{ViewManager}, and UI components managed via callbacks. \\ \hline
    \end{tabular}
    \label{tab:CEP}
\end{table}

\section{Design Goals and Objectives}
The main objectives of this project are:
\begin{itemize}
    \item To implement a fully functional local music player using Java and JavaFX.
    \item To apply object-oriented design principles including encapsulation, inheritance, and polymorphism.
    \item To demonstrate the MVC architectural pattern in a real-world application.
    \item To provide a modern, responsive UI that adapts to different window sizes.
    \item To ensure data persistence so the library survives application restarts.
    \item To support background processing for metadata loading without blocking the UI.
\end{itemize}

\section{Application Scope}
Auralin Music Player is intended for personal desktop use on Windows. It supports MP3, WAV, FLAC, and M4A audio formats for import, with full playback support for MP3 files via the JavaFX \texttt{MediaPlayer} API. The application can be packaged as a standalone JAR or installed via an Inno Setup installer with Windows shell integration.


% ============================================================
% CHAPTER 2: SYSTEM ARCHITECTURE
% ============================================================
\newpage
\chapter{System Architecture}

\section{Architectural Overview}
Auralin follows a layered \textbf{Model-View-Controller (MVC)} architecture. The codebase is organised into the following packages, each with a clearly defined responsibility:

\begin{table}[htbp]
\centering
\caption{Package structure and responsibilities}
\begin{tabular}{|p{4cm}|p{10cm}|}
\hline
\textbf{Package} & \textbf{Responsibility} \\ \hline
\texttt{model} & Data classes (\texttt{Song}) representing domain entities. \\ \hline
\texttt{services} & \texttt{DatabaseService} --- library management, persistence, metadata loading. \\ \hline
\texttt{controllers} & \texttt{PlayerController} --- playback logic, MediaPlayer lifecycle. \\ \hline
\texttt{managers} & \texttt{ViewManager} --- view switching, table refresh, scroll-bar styling. \\ \hline
\texttt{views} & \texttt{HomeView}, \texttt{LibraryView}, \texttt{SearchView}, \texttt{LikedView}, \texttt{NowPlayingView} --- screen layouts. \\ \hline
\texttt{components} & \texttt{PlayerBar}, \texttt{Sidebar}, \texttt{MusicCard} --- reusable UI widgets. \\ \hline
\texttt{ui} & \texttt{SongTable} --- custom \texttt{TableView} with artwork, like toggle, and row highlighting. \\ \hline
\texttt{utils} & \texttt{FileImportService} --- file/folder chooser dialogs and background scanning. \\ \hline
\texttt{config} & \texttt{Constants} --- colour palette, CSS strings, default artwork path. \\ \hline
\end{tabular}
\end{table}

\section{Class Diagram Description}
The following describes the key relationships between the main classes:

\begin{itemize}
    \item \textbf{\texttt{Main}} (Application entry point) creates and wires together \texttt{DatabaseService}, \texttt{PlayerController}, \texttt{ViewManager}, \texttt{Sidebar}, \texttt{PlayerBar}, and \texttt{FileImportService}. It owns the root \texttt{StackPane} and the \texttt{NowPlayingView} overlay.

    \item \textbf{\texttt{DatabaseService}} owns an \texttt{ObservableList<Song>} and an \texttt{ExecutorService} thread pool. It exposes methods to add files, retrieve all/liked/recent/recommended songs, record plays, and persist/load the library.

    \item \textbf{\texttt{PlayerController}} holds a reference to \texttt{DatabaseService} and \texttt{PlayerBar}. It manages the \texttt{MediaPlayer} lifecycle and fires three callbacks: \texttt{onSongChange}, \texttt{onLikeChange}, and \texttt{onPlayStateChange}.

    \item \textbf{\texttt{ViewManager}} holds references to \texttt{DatabaseService}, the content \texttt{StackPane}, and \texttt{PlayerController}. It creates view instances on demand and manages the \texttt{ScrollPane} wrapper.

    \item \textbf{\texttt{Song}} is a plain data class with title, artist, file path, artwork (\texttt{Image}), duration, play count, last-played timestamp, and liked flag.

    \item \textbf{\texttt{PlayerBar}} is a composite \texttt{HBox} containing artwork, song info, playback controls (SVG play/pause button, shuffle, repeat), a progress slider, and a volume slider. It communicates upward via the \texttt{PlayerBarListener} interface.

    \item \textbf{\texttt{Sidebar}} is a \texttt{VBox} with navigation buttons (Home, Search, Library, Liked Songs) and a loading progress bar. It communicates via \texttt{SidebarListener}.

    \item \textbf{\texttt{SongTable}} extends \texttt{TableView<Song>} and renders title/artwork, artist, duration, and a heart like-toggle column.

    \item \textbf{\texttt{MusicCard}} is a \texttt{VBox} card used in the Home view grid, showing album art, title, artist, and a play/pause overlay for the currently playing song.
\end{itemize}

\section{Data Flow}
The data flow through the application follows this sequence:

\begin{enumerate}
    \item The user imports files via \texttt{FileImportService}, which calls \texttt{DatabaseService.addFile()}.
    \item \texttt{DatabaseService} creates a \texttt{Song} object and submits a background task to load ID3 metadata via \texttt{mp3agic}.
    \item When metadata is ready, \texttt{Platform.runLater()} updates the \texttt{ObservableList<Song>} on the JavaFX thread.
    \item The user double-clicks a song in \texttt{SongTable} or clicks a \texttt{MusicCard}, triggering \texttt{PlayerController.play(Song)}.
    \item \texttt{PlayerController} creates a \texttt{MediaPlayer}, registers time/end listeners, and fires \texttt{onSongChange} to refresh views.
    \item \texttt{ViewManager} receives the callback and refreshes the active view or just the table rows, depending on which view is active.
\end{enumerate}

\section{Technology Stack}
\begin{table}[htbp]
\centering
\caption{Technology stack}
\begin{tabular}{|p{4cm}|p{10cm}|}
\hline
\textbf{Technology} & \textbf{Usage} \\ \hline
Java 17 & Core application language \\ \hline
JavaFX 17 & UI framework (Scene Graph, CSS, MediaPlayer) \\ \hline
mp3agic & ID3v1/v2 tag parsing and duration extraction \\ \hline
Inno Setup & Windows installer packaging \\ \hline
\end{tabular}
\end{table}


% ============================================================
% CHAPTER 3: FEATURES AND FUNCTIONALITY
% ============================================================
\newpage
\chapter{Features and Functionality}

\section{Library Management}
\begin{itemize}
    \item \textbf{Add Songs:} Users can import individual audio files (MP3, WAV, FLAC, M4A) via a file chooser dialog.
    \item \textbf{Add Folder:} Users can import an entire directory recursively. A progress bar in the sidebar shows loading status (e.g., ``Loading: 12/45'').
    \item \textbf{Persistent Library:} The library is saved to \texttt{\%APPDATA\%\textbackslash AuralinPlayer\textbackslash library.dat} and restored on next launch. Songs that no longer exist on disk are silently skipped.
    \item \textbf{Duplicate Prevention:} \texttt{DatabaseService.addFile()} checks whether a file path already exists in the library before adding.
\end{itemize}

\section{Playback Controls}
\begin{itemize}
    \item \textbf{Play / Pause:} A premium SVG play/pause button (white circle with black icon) toggles playback. The icon swaps between a triangle (play) and two bars (pause).
    \item \textbf{Next / Previous:} Advances or retreats through the song list, with wrap-around at both ends.
    \item \textbf{Shuffle:} Picks a random song from the library on each ``next'' action. Mutually exclusive with Repeat.
    \item \textbf{Repeat:} Restarts the current song when it ends. Mutually exclusive with Shuffle.
    \item \textbf{Seek:} A progress slider allows scrubbing to any position in the track.
    \item \textbf{Volume:} A volume slider controls playback volume (0.0--1.0).
    \item \textbf{Auto-advance:} When a song ends and Repeat is off, the player automatically advances to the next song.
\end{itemize}

\section{Liked Songs}
\begin{itemize}
    \item Any song can be liked/unliked via the heart button in the \texttt{PlayerBar}, the heart column in \texttt{SongTable}, or the heart button in \texttt{NowPlayingView}.
    \item Liked status is persisted to disk immediately on toggle.
    \item The \textbf{Liked Songs} view shows only liked songs using a filtered \texttt{ObservableList}.
    \item When a song is unliked from the Liked view, it disappears from the list instantly.
\end{itemize}

\section{Home Screen}
\begin{itemize}
    \item \textbf{Recently Played:} Shows up to 9 songs sorted by \texttt{lastPlayed} timestamp descending.
    \item \textbf{Made For You:} Shows up to 12 songs sorted by play count descending (most-played recommendations).
    \item Both sections use a responsive \texttt{FlowPane} grid of \texttt{MusicCard} components that reflows on window resize.
    \item The currently playing song shows a red play/pause overlay on its card.
\end{itemize}

\section{Search}
The Search view displays all songs in a \texttt{SongTable} with a search field that filters by title or artist in real time using JavaFX \texttt{FilteredList}.

\section{Now Playing Overlay}
Clicking the album artwork in the \texttt{PlayerBar} opens a full-screen \texttt{NowPlayingView} overlay that covers the entire window (including sidebar and player bar). It shows:
\begin{itemize}
    \item Large centered album artwork with a drop shadow.
    \item Song title and artist name.
    \item A heart button to toggle liked status without closing the overlay.
    \item A close button (top-right) and Escape key support.
\end{itemize}

\section{Open With Integration}
The application registers itself as a handler for \texttt{.mp3} files via a Windows registry entry (\texttt{register-open-with.reg}). When launched with a file argument, it adds the file to the library and begins playback after a short delay to allow the \texttt{MediaPlayer} to initialise.

\section{Custom Window Chrome}
\begin{itemize}
    \item The application uses \texttt{StageStyle.TRANSPARENT} with a custom title bar, eliminating the native OS window frame.
    \item The window has rounded corners (10px radius) and a subtle dark border.
    \item Custom minimize, maximize, and close buttons with Windows 11-style hover animations (fade and scale transitions).
    \item All 8 edge/corner resize handles are implemented via transparent overlay panes.
    \item The title bar supports drag-to-move and double-click-to-maximize.
\end{itemize}


% ============================================================
% CHAPTER 4: IMPLEMENTATION DETAILS
% ============================================================
\newpage
\chapter{Implementation Details}

\section{Key Classes}

\subsection{Main.java}
\texttt{Main} extends \texttt{javafx.application.Application} and serves as the composition root. It:
\begin{itemize}
    \item Instantiates all major components and wires them together.
    \item Registers three callbacks on \texttt{PlayerController}: \texttt{onSongChange}, \texttt{onLikeChange}, and \texttt{onPlayStateChange}.
    \item Builds the outer layout: a \texttt{BorderPane} with a custom title bar at top, a \texttt{BorderPane} (sidebar + content area) in the centre, and \texttt{PlayerBar} at the bottom.
    \item Wraps everything in a \texttt{StackPane} for the \texttt{NowPlayingView} overlay.
    \item Sets up 8 transparent resize panes anchored to the window edges.
    \item Handles the ``Open With'' file argument.
    \item Writes global CSS to a temp file and loads it into the scene.
\end{itemize}

\subsection{DatabaseService.java}
\texttt{DatabaseService} is the application's data layer. Key implementation details:
\begin{itemize}
    \item Maintains an \texttt{ObservableList<Song>} as the single source of truth for the library.
    \item Uses a fixed thread pool of 3 threads (\texttt{Executors.newFixedThreadPool(3)}) for background metadata loading and disk I/O.
    \item Metadata is loaded via \texttt{mp3agic}'s \texttt{Mp3File} class, which reads ID3v1/v2 tags without creating a \texttt{MediaPlayer}.
    \item The save format is a pipe-delimited text file: \texttt{filePath|liked|plays|lastPlayed|artist|durationMs}.
    \item On load, songs whose files no longer exist are silently skipped.
    \item \texttt{getRecentlyPlayed()} and \texttt{getRecommendations()} use Java Stream API with \texttt{Comparator} and \texttt{limit()}.
\end{itemize}

\subsection{PlayerController.java}
\texttt{PlayerController} manages the \texttt{javafx.scene.media.MediaPlayer} lifecycle:
\begin{itemize}
    \item \texttt{play(Song)} disposes the previous \texttt{MediaPlayer}, creates a new one, sets volume, registers time/end listeners, and starts playback.
    \item \texttt{togglePlay()} checks \texttt{MediaPlayer.Status.PLAYING} to decide whether to pause or resume.
    \item \texttt{playNext()} supports both sequential and shuffle modes.
    \item \texttt{setRepeatMode(1)} causes \texttt{setOnEndOfMedia} to seek to \texttt{Duration.ZERO} and replay.
    \item Three \texttt{Runnable} callbacks (\texttt{onSongChange}, \texttt{onLikeChange}, \texttt{onPlayStateChange}) are fired via \texttt{Platform.runLater()} to ensure UI updates happen on the JavaFX Application Thread.
\end{itemize}

\subsection{ViewManager.java}
\texttt{ViewManager} acts as a mediator between the navigation system and the view layer:
\begin{itemize}
    \item \texttt{showView(String)} creates a fresh view instance, wraps it in a \texttt{ScrollPane}, and replaces the content area.
    \item \texttt{refreshTableOnly()} walks the scene graph to find the \texttt{SongTable} and calls \texttt{refresh()} on it, avoiding full view recreation (which would cause header flicker).
    \item \texttt{refreshHomeCards()} only rebuilds the Home view if it is currently active.
    \item Scroll bar nodes are styled to be invisible after each layout pass via \texttt{Platform.runLater()}.
\end{itemize}

\subsection{Song.java}
\texttt{Song} is a plain Java object (POJO) with:
\begin{itemize}
    \item \texttt{title}, \texttt{artist}, \texttt{filePath} (String fields).
    \item \texttt{artwork} (\texttt{javafx.scene.image.Image}).
    \item \texttt{duration} (\texttt{javafx.util.Duration}).
    \item \texttt{plays} (int) and \texttt{lastPlayed} (long, epoch milliseconds) --- public fields for simplicity.
    \item \texttt{liked} (boolean).
    \item \texttt{record()} increments \texttt{plays} and sets \texttt{lastPlayed} to \texttt{System.currentTimeMillis()}.
    \item \texttt{getDurationStr()} formats duration as \texttt{m:ss}.
\end{itemize}

\section{Design Patterns Used}

\subsection{Model-View-Controller (MVC)}
The application strictly separates:
\begin{itemize}
    \item \textbf{Model:} \texttt{Song}, \texttt{DatabaseService} (data and business logic).
    \item \textbf{View:} All classes in \texttt{views}, \texttt{components}, and \texttt{ui} packages (JavaFX nodes).
    \item \textbf{Controller:} \texttt{PlayerController} (playback logic), \texttt{ViewManager} (navigation logic).
\end{itemize}

\subsection{Observer Pattern}
\texttt{PlayerController} exposes three \texttt{Runnable} callback slots (\texttt{setOnSongChange}, \texttt{setOnLikeChange}, \texttt{setOnPlayStateChange}). \texttt{Main} registers lambdas that call \texttt{ViewManager} methods, effectively implementing the Observer pattern without a formal event bus.

\subsection{Strategy Pattern}
The \texttt{PlayerBarListener}, \texttt{SidebarListener}, \texttt{SongTableListener}, \texttt{MusicCardListener}, and \texttt{NowPlayingListener} interfaces define strategies for handling user interactions. Each is implemented as an anonymous class or lambda at the call site, allowing the UI components to remain decoupled from business logic.

\subsection{Facade Pattern}
\texttt{DatabaseService} acts as a facade over the underlying \texttt{ObservableList}, file I/O, and thread pool. Callers interact with a simple API (\texttt{addFile}, \texttt{getAllSongs}, \texttt{saveToDisk}, etc.) without knowing the implementation details.

\subsection{Template Method (Implicit)}
\texttt{Main.start(Stage)} follows the JavaFX application lifecycle template: \texttt{init()} $\rightarrow$ \texttt{start()} $\rightarrow$ \texttt{stop()}. The \texttt{stop()} override disposes the player and saves the library.

\section{Concurrency Model}
All metadata loading and disk I/O run on a background \texttt{ExecutorService} thread pool. UI updates are always marshalled back to the JavaFX Application Thread via \texttt{Platform.runLater()}. This prevents UI freezing during large folder imports.

\section{CSS Styling}
Global CSS is defined as a Java string in \texttt{Constants.GLOBAL\_CSS} and written to a temporary file at startup. This approach avoids issues with data URI encoding of special characters. The CSS defines styles for \texttt{.nav-btn}, \texttt{.action-btn}, \texttt{.music-card}, \texttt{.spotify-table}, \texttt{.spotify-slider}, \texttt{.control-icon}, and \texttt{.search-field}.


% ============================================================
% CHAPTER 5: DATABASE / PERSISTENCE DESIGN
% ============================================================
\newpage
\chapter{Database and Persistence Design}

\section{Persistence Strategy}
Auralin does not use a relational database. Instead, it uses a simple flat-file persistence model: a pipe-delimited text file stored in the user's application data directory.

\subsection{Save File Location}
\begin{itemize}
    \item \textbf{Windows:} \texttt{\%APPDATA\%\textbackslash AuralinPlayer\textbackslash library.dat}
    \item \textbf{Other OS:} \texttt{\textasciitilde/.auralin/library.dat}
\end{itemize}
The directory is created automatically if it does not exist.

\subsection{File Format}
Each line in \texttt{library.dat} represents one song with six pipe-separated fields:

\begin{center}
\texttt{filePath|liked|plays|lastPlayed|artist|durationMs}
\end{center}

\begin{table}[htbp]
\centering
\caption{Library file fields}
\begin{tabular}{|p{3cm}|p{3cm}|p{7cm}|}
\hline
\textbf{Field} & \textbf{Type} & \textbf{Description} \\ \hline
\texttt{filePath} & String & Absolute path to the audio file \\ \hline
\texttt{liked} & boolean & Whether the song is liked \\ \hline
\texttt{plays} & int & Number of times the song has been played \\ \hline
\texttt{lastPlayed} & long & Unix epoch milliseconds of last play \\ \hline
\texttt{artist} & String & Artist name from ID3 tag \\ \hline
\texttt{durationMs} & long & Track duration in milliseconds \\ \hline
\end{tabular}
\end{table}

\subsection{Save Trigger}
The library is saved:
\begin{itemize}
    \item After every \texttt{addFile()} call (background thread).
    \item After every \texttt{recordPlay()} call (when a song starts playing).
    \item After every like toggle (\texttt{toggleLike()} in \texttt{PlayerController} and \texttt{onLikeToggle()} in \texttt{SongTable}).
    \item On application shutdown (\texttt{Main.stop()}).
\end{itemize}

\subsection{Load Process}
On startup, \texttt{DatabaseService} reads \texttt{library.dat} line by line. For each line:
\begin{enumerate}
    \item It checks whether the file still exists on disk; if not, the entry is skipped.
    \item It creates a \texttt{Song} object and populates all fields from the saved data.
    \item It submits a background task to load the album artwork via \texttt{mp3agic} (artwork is not stored in the file).
\end{enumerate}

\section{Metadata Loading}
Album artwork and artist names are extracted from ID3 tags using the \textbf{mp3agic} library. This library reads the MP3 file header directly without creating a \texttt{MediaPlayer}, making it significantly faster and more memory-efficient. The process runs on a background thread and updates the \texttt{ObservableList} via \texttt{Platform.runLater()} when complete.

\section{Filtered Views}
\begin{itemize}
    \item \texttt{getLikedSongs()} returns \texttt{all.filtered(Song::isLiked)}, a live \texttt{FilteredList} that automatically updates when the liked flag changes.
    \item \texttt{getRecentlyPlayed()} and \texttt{getRecommendations()} use Java Streams and return plain \texttt{List<Song>} snapshots (not live lists), as the Home view is rebuilt on each navigation.
\end{itemize}


% ============================================================
% CHAPTER 6: UI DESIGN
% ============================================================
\newpage
\chapter{UI Design}

\section{Design Philosophy}
The UI is inspired by Spotify's dark interface. The colour palette is defined in \texttt{Constants.java} and uses:
\begin{itemize}
    \item \textbf{Background:} \texttt{\#121212} (dark grey) and \texttt{\#000000} (pure black).
    \item \textbf{Sidebar:} \texttt{\#000000}.
    \item \textbf{Cards:} \texttt{\#181818} with \texttt{\#2A2A2A} on hover.
    \item \textbf{Accent:} \texttt{\#FA2D48} (red) --- used for the logo, active nav items, liked hearts, and the current song highlight.
    \item \textbf{Secondary text:} \texttt{\#B3B3B3} (grey).
\end{itemize}

\section{Layout Structure}
The main window is a \texttt{BorderPane} with:
\begin{itemize}
    \item \textbf{Top:} Custom title bar (32px height) with app icon, title label, and window control buttons.
    \item \textbf{Left:} \texttt{Sidebar} (240px wide) with navigation buttons and loading progress bar.
    \item \textbf{Centre:} \texttt{StackPane} (content area) wrapped in a \texttt{ScrollPane}, showing the active view.
    \item \textbf{Bottom:} \texttt{PlayerBar} (120px height) with now-playing info, controls, and sliders.
\end{itemize}
A \texttt{NowPlayingView} overlay can be pushed onto the root \texttt{StackPane} to cover the entire window.

\section{Sidebar Component}
The \texttt{Sidebar} contains:
\begin{itemize}
    \item A red ``Auralin'' logo label (28px bold).
    \item Four navigation buttons: Home (house emoji), Search (magnifier emoji), Your Library (books emoji), and Liked Songs (SVG heart icon).
    \item Active navigation item is highlighted in red; inactive items are grey with a white hover effect.
    \item A \texttt{ProgressBar} and status label that appear during folder import.
\end{itemize}
The Liked Songs button uses a custom SVG heart icon (loaded from \texttt{/resources/heart.svg}) instead of an emoji, for visual consistency with the heart icons in \texttt{PlayerBar} and \texttt{SongTable}.

\section{Player Bar Component}
The \texttt{PlayerBar} is divided into three sections:
\begin{itemize}
    \item \textbf{Left (info):} 60$\times$60px album artwork (clickable to open Now Playing), song title, artist name, and heart like button.
    \item \textbf{Centre (controls):} Shuffle, Previous, Play/Pause (white circle with SVG icon), Next, Repeat buttons, and a progress slider with current/total time labels.
    \item \textbf{Right (volume):} Volume icon and volume slider.
\end{itemize}
The progress and volume sliders use a custom CSS class \texttt{.spotify-slider} with a red-filled track that updates dynamically via inline style injection.

\section{Song Table Component}
\texttt{SongTable} extends \texttt{TableView<Song>} with four columns:
\begin{itemize}
    \item \textbf{TITLE:} Shows 40$\times$40px artwork thumbnail, song title (red if currently playing), and artist name.
    \item \textbf{ARTIST:} Plain text artist name.
    \item \textbf{DURATION:} Right-aligned formatted duration string.
    \item \textbf{LIKED:} SVG heart toggle button with pop animation on click.
\end{itemize}
The table uses \texttt{CONSTRAINED\_RESIZE\_POLICY} and hides all scroll bars. Row hover is styled via CSS (\texttt{.spotify-table .table-row-cell:hover}).

\section{Music Card Component}
\texttt{MusicCard} is a 160$\times$\textasciitilde{}220px \texttt{VBox} card used in the Home view grid:
\begin{itemize}
    \item Square album artwork (140$\times$140px) with rounded corners (12px radius).
    \item A red circle play/pause overlay at the bottom-right corner of the artwork, visible only for the currently playing song.
    \item Song title (bold, red if current) and artist name below the artwork.
    \item Smooth 1.05$\times$ scale-up animation on hover (150ms).
\end{itemize}

\section{Now Playing View}
\texttt{NowPlayingView} is a full-screen \texttt{StackPane} overlay with:
\begin{itemize}
    \item Pure black background.
    \item ``Now Playing'' heading (top-centre, 20px bold white).
    \item Close button (top-right, grey ``$\times$'' that turns white on hover).
    \item 340$\times$340px album artwork with drop shadow.
    \item Song title (26px bold white) and artist (16px grey) centred below the artwork.
    \item Heart like button with pop animation and red glow effect when liked.
    \item Escape key closes the overlay.
\end{itemize}

\section{Responsive Design}
\begin{itemize}
    \item The window size is capped at 1250$\times$850px but adapts to smaller screens using \texttt{Screen.getPrimary().getVisualBounds()}.
    \item The Home view \texttt{FlowPane} binds its \texttt{prefWrapLength} to the view width, so cards reflow correctly on window resize.
    \item Minimum window size is 600$\times$400px.
\end{itemize}


% ============================================================
% CHAPTER 7: CHALLENGES AND SOLUTIONS
% ============================================================
\newpage
\chapter{Challenges and Solutions}

\section{UI Thread Safety}
\textbf{Challenge:} JavaFX requires all UI updates to occur on the JavaFX Application Thread. Background metadata loading and disk I/O run on worker threads, so updating the \texttt{ObservableList} or UI nodes directly from those threads causes \texttt{IllegalStateException}.

\textbf{Solution:} All UI updates from background threads are wrapped in \texttt{Platform.runLater(Runnable)}. The \texttt{DatabaseService} executor submits tasks that compute results on worker threads and then call \texttt{Platform.runLater()} to apply changes to the \texttt{ObservableList}.

\section{Table Header Flickering on Like Toggle}
\textbf{Challenge:} When a user toggles the like button in \texttt{SongTable}, the original approach was to call \texttt{ViewManager.refreshCurrentView()}, which recreated the entire view including the \texttt{TableView}. This caused the column headers to flicker visibly.

\textbf{Solution:} A dedicated \texttt{refreshTableOnly()} method was added to \texttt{ViewManager}. It walks the scene graph to find the existing \texttt{SongTable} instance and calls \texttt{TableView.refresh()} on it, which re-renders only the cell content without recreating the headers.

\section{Artwork Not Loading on Startup}
\textbf{Challenge:} On application startup, artwork was not displayed for songs loaded from \texttt{library.dat} because the original implementation relied on \texttt{MediaPlayer.setOnReady()} to extract metadata, which required creating a \texttt{MediaPlayer} for every song --- an expensive operation.

\textbf{Solution:} The \texttt{mp3agic} library was integrated to read ID3 tags directly from the MP3 file header without creating a \texttt{MediaPlayer}. This is both faster and more memory-efficient. Artwork loading is submitted as a background task and applied via \texttt{Platform.runLater()}.

\section{Shuffle and Repeat Mutual Exclusion}
\textbf{Challenge:} Having both shuffle and repeat active simultaneously leads to undefined behaviour (shuffle picks a random next song, but repeat restarts the current song --- they conflict).

\textbf{Solution:} The \texttt{PlayerBar} buttons implement mutual exclusion: activating shuffle programmatically turns off repeat (and vice versa), updating both the button styles and firing the appropriate listener callbacks.

\section{Custom Window Resize}
\textbf{Challenge:} Using \texttt{StageStyle.TRANSPARENT} removes the native OS window frame, including resize handles. Implementing resize from scratch requires tracking mouse press/drag events on all 8 edges and corners.

\textbf{Solution:} Eight transparent \texttt{Pane} nodes are anchored to the window edges and corners using \texttt{AnchorPane}. Each pane sets the appropriate \texttt{Cursor} and handles \texttt{setOnMousePressed}/\texttt{setOnMouseDragged} events to compute new window position and size, respecting the minimum window dimensions.

\section{Open With Shell Integration}
\textbf{Challenge:} When Windows launches the application with a file path argument (via ``Open With''), the \texttt{MediaPlayer} may not be ready immediately after \texttt{addFile()} is called, causing the song to start with unknown duration.

\textbf{Solution:} A \texttt{PauseTransition} of 800ms is inserted between adding the file and calling \texttt{PlayerController.play()}, giving the \texttt{MediaPlayer} time to initialise and report the correct duration.

\section{CSS Loading with Special Characters}
\textbf{Challenge:} Loading CSS via a \texttt{data:text/css,} URI requires URL-encoding special characters (\texttt{\#}, spaces, parentheses, etc.), which is error-prone and hard to maintain.

\textbf{Solution:} The CSS string is written to a temporary file at startup (\texttt{File.createTempFile()}), and the file's URI is loaded into the scene's stylesheet list. The temp file is marked for deletion on JVM exit.


% ============================================================
% CHAPTER 8: CONCLUSION
% ============================================================
\newpage
\chapter{Conclusion}

Auralin Music Player successfully demonstrates the application of object-oriented programming principles and the MVC architectural pattern in a real-world JavaFX desktop application. The project achieves its primary objectives: a visually polished, Spotify-inspired music player that manages a local MP3 library, persists data between sessions, and provides a smooth, responsive user experience.

\section{Key Achievements}
\begin{itemize}
    \item \textbf{Clean Architecture:} The codebase is well-organised into packages with clear separation of concerns. The \texttt{model}, \texttt{services}, \texttt{controllers}, \texttt{managers}, \texttt{views}, \texttt{components}, \texttt{ui}, \texttt{utils}, and \texttt{config} packages each have a single, well-defined responsibility.
    \item \textbf{Design Patterns:} MVC, Observer, Strategy, and Facade patterns are applied consistently throughout the codebase.
    \item \textbf{Performance:} Background metadata loading via \texttt{mp3agic} and a thread pool ensures the UI remains responsive even when importing large music folders.
    \item \textbf{Persistence:} A lightweight flat-file persistence mechanism saves and restores the complete library state, including play counts, liked status, and last-played timestamps.
    \item \textbf{UI Quality:} The custom frameless window, SVG icons, smooth animations, and responsive layout deliver a premium user experience comparable to commercial music players.
    \item \textbf{Extensibility:} The listener/callback interfaces make it straightforward to add new views, playback features, or data sources without modifying existing classes.
\end{itemize}

\section{Limitations}
\begin{itemize}
    \item The application currently supports only local files; streaming or network sources are not supported.
    \item Playlist management (creating, editing, and saving custom playlists) is not yet implemented.
    \item The search functionality filters only by title and artist; album and genre search are not supported.
    \item The flat-file persistence format does not support concurrent access from multiple instances.
\end{itemize}

\section{Future Work}
\begin{itemize}
    \item \textbf{Playlist support:} Allow users to create, name, and manage custom playlists.
    \item \textbf{Equaliser:} Integrate a graphic equaliser using JavaFX \texttt{AudioEqualizer}.
    \item \textbf{Lyrics display:} Fetch and display song lyrics from an online API.
    \item \textbf{Cross-platform packaging:} Use \texttt{jpackage} to create native installers for macOS and Linux.
    \item \textbf{Database migration:} Replace the flat-file store with SQLite for better query performance and concurrent access support.
    \item \textbf{Album view:} Group songs by album and display a dedicated album grid view.
\end{itemize}

In summary, Auralin Music Player is a well-engineered, feature-complete music player that demonstrates strong object-oriented design, effective use of JavaFX APIs, and thoughtful UI/UX decisions. It serves as a solid foundation for further development into a full-featured desktop music application.

% ============================================================
% REFERENCES
% ============================================================
\newpage
\chapter*{References}
\addcontentsline{toc}{chapter}{References}

\begin{thebibliography}{99}

\bibitem{javafx}
Oracle Corporation,
\textit{JavaFX Documentation},
Available: \url{https://openjfx.io/javadoc/17/},
Accessed: December 2025.

\bibitem{mp3agic}
M. Karger,
\textit{mp3agic: A Java library for reading and manipulating MP3 files},
Available: \url{https://github.com/mpatric/mp3agic},
Accessed: December 2025.

\bibitem{javamedia}
Oracle Corporation,
\textit{JavaFX Media: MediaPlayer API},
Available: \url{https://openjfx.io/javadoc/17/javafx.media/javafx/scene/media/MediaPlayer.html},
Accessed: December 2025.

\bibitem{mvc}
E. Gamma, R. Helm, R. Johnson, and J. Vlissides,
\textit{Design Patterns: Elements of Reusable Object-Oriented Software},
Addison-Wesley, 1994.

\bibitem{javaconcurrency}
B. Goetz, T. Peierls, J. Bloch, J. Bowbeer, D. Holmes, and D. Lea,
\textit{Java Concurrency in Practice},
Addison-Wesley, 2006.

\bibitem{effectivejava}
J. Bloch,
\textit{Effective Java}, 3rd ed.,
Addison-Wesley, 2018.

\bibitem{javafxcss}
Oracle Corporation,
\textit{JavaFX CSS Reference Guide},
Available: \url{https://openjfx.io/javadoc/17/javafx.graphics/javafx/scene/doc-files/cssref.html},
Accessed: December 2025.

\end{thebibliography}

\end{document}
